% SGA 3, Expose IV, Proposition 4.5.6 through Remark 4.5.9.
% French authority: Exp4-13oct24.pdf, local pages 46--47
% (combined-reader pages 234--235; printed pages 222--223).
% The external English Markdown was used only as a comparison.

\medskip
\noindent\SGARefTarget{sga3:iv:proposition:4:5:6}\textbf{命题 4.5.6.}
\textit{设 \(F\) 和 \(G\) 是 \(S\) 上的两个层,
\(f:F\to G\) 是一个 \(S\)-态射.
关于 \(f\) 的以下条件等价:}
\begin{enumerate}
  \item[(i)]\SGARefTarget{sga3:iv:proposition:4:5:6:item:i:01}
  \textit{\(f\) 是 \(\widetilde{\mathscr C}\) 中的单态射
  (分别地, 满态射; 分别地, 同构);}
  \item[(ii)]\SGARefTarget{sga3:iv:proposition:4:5:6:item:ii:01}
  \textit{\(f\) 是 \(\widetilde{\mathscr C}_{/S}=\widetilde{\mathscr C/S}\)
  中的单态射 (分别地, 满态射; 分别地, 同构).}
\end{enumerate}

对单态射与同构, 这是直接的, 因为所讨论的是关于预层的问题.
对满态射, 结论由将满态射描述为覆盖态射的结果
(\SGARefLink{sga3:iv:proposition:4:4:3}{4.4.3}) 得到,
因为按诱导拓扑的\SGARefLink{sga3:iv:subsection:4:5}{定义},
覆盖态射在 \(\widetilde{\mathscr C}\) 中与在
\(\widetilde{\mathscr C/S}\) 中是相同的.

\medskip
\noindent\SGARefTarget{sga3:iv:proposition:4:5:7}\textbf{命题 4.5.7.}
\textit{设 \(f:F\to G\) 是一个层态射. 以下条件等价:}
\begin{enumerate}
  \item[(i)]\SGARefTarget{sga3:iv:proposition:4:5:7:item:i:01}
  \textit{\(f\) 是单态射 (分别地, 满态射; 分别地, 同构);}
  \item[(ii)]\SGARefTarget{sga3:iv:proposition:4:5:7:item:ii:01}
  \textit{对每个 \(S\in\operatorname{Ob}(\mathscr C)\), 态射
  \(f_S:F_S\to G_S\) 局部为单态射
  (分别地, 满态射; 分别地, 同构), 即:}
  \item[(iii)]\SGARefTarget{sga3:iv:proposition:4:5:7:item:iii:01}
  \textit{对每个 \(S\in\operatorname{Ob}(\mathscr C)\),
  使 \(F_T\to G_T\) 为单态射 (分别地, 满态射; 分别地, 同构)
  的那些态射 \(T\to S\) 所成的集合, 是 \(S\) 的一个加细.}
\end{enumerate}
\textit{若给定的拓扑由预拓扑 \(\mathscr R\) 定义,
则这些条件还等价于:}
\begin{enumerate}
  \item[(iv)]\SGARefTarget{sga3:iv:proposition:4:5:7:item:iv:01}
  \textit{对每个 \(S\in\operatorname{Ob}(\mathscr C)\), 存在一个覆盖族
  \(\{S_i\to S\}\in\mathscr R(S)\),
  使得对每个 \(i\), \(F_{S_i}\to G_{S_i}\) 都是单态射
  (分别地, 满态射; 分别地, 同构).}
\end{enumerate}
\textit{若 \(\mathscr C\) 有终对象 \(e\),
则在条件 \SGARefLink{sga3:iv:proposition:4:5:7:item:ii:01}{(ii)},
\SGARefLink{sga3:iv:proposition:4:5:7:item:iii:01}{(iii)} 和
\SGARefLink{sga3:iv:proposition:4:5:7:item:iv:01}{(iv)} 中只需取 \(S=e\).}

显然, \SGARefLink{sga3:iv:proposition:4:5:7:item:ii:01}{(ii)},
\SGARefLink{sga3:iv:proposition:4:5:7:item:iii:01}{(iii)} 和
\SGARefLink{sga3:iv:proposition:4:5:7:item:iv:01}{(iv)} 等价.
要证明 \SGARefLink{sga3:iv:proposition:4:5:7:item:i:01}{(i)} 与
\SGARefLink{sga3:iv:proposition:4:5:7:item:ii:01}{(ii)} 等价,
以及关于终对象的\SGARefLink{sga3:iv:proposition:4:5:7}{断言},
只需证明 \SGARefLink{sga3:iv:proposition:4:5:7:item:ii:01}{(ii)} 蕴含
\SGARefLink{sga3:iv:proposition:4:5:7:item:i:01}{(i)},
且所讨论的性质在基变换下稳定.
先证明后一事实. 对单态射与同构, 这是直接的,
因为所讨论的是关于预层的问题.
对满态射, 则由每个层满态射都是普遍的这一事实得到
(\SGARefLink{sga3:iv:proposition:4:4:3}{4.4.3}).

还需证明 \SGARefLink{sga3:iv:proposition:4:5:7:item:ii:01}{(ii)} 蕴含
\SGARefLink{sga3:iv:proposition:4:5:7:item:i:01}{(i)}.
先假设 \(f_S:F_S\to G_S\) 局部为单态射 (分别地, 同构).
此时存在 \(S\) 的一个覆盖筛 \(C\),
使得对属于 \(C\) 的每个 \(T\to S\),
\(f_T\) 都是单态射 (分别地, 同构).
由于单态射 (分别地, 同构) 的逆极限仍是单态射 (分别地, 同构),
\(\operatorname{Hom}(C,f)\) 是单态射 (分别地, 同构);
见 \SGARefLink{sga3:iv:proposition:4:1:4}{4.1.4}.
所以交换图
\[
\begin{array}{ccc}
 \operatorname{Hom}(C,F)
 &\xrightarrow{\operatorname{Hom}(C,f)}&
 \operatorname{Hom}(C,G)
 \\
 \big\uparrow\wr && \big\uparrow\wr
 \\
 F(S)&\xrightarrow{f(S)}&G(S)
\end{array}
\]
表明 \(f(S)\) 是单射 (分别地, 双射).

最后, 假设 \(f:F\to G\) 局部为满态射,
设 \(H\subset G\) 是它的像.
对每个 \(S\to G\), 层 \(H\times_GS\) 是
\[
  f\times_GS:F\times_GS\longrightarrow S.
\]
的像. 要证明 \(f\) 是满态射, 只需证明 \(H\) 是 \(G\) 的一个加细,
或等价地, 对每个 \(S\), \(H\times_GS\) 是 \(S\) 的一个加细.
但由于 \(f\times_GS\) 局部为覆盖态射,
这一点在沿属于 \(S\) 的一个加细的每个 \(T\to S\) 作基变换后成立.
所以, 由公理 \SGARefLink{sga3:iv:definition:4:2:1:item:t:2:01}{(T~2)},
\(H\times_GS\) 本身就是 \(S\) 的一个加细.

\medskip
\noindent\SGARefTarget{sga3:iv:corollary:4:5:8}\textbf{推论 4.5.8.}
\textit{设 \(F\) 与 \(G\) 是 \(S\) 上的两个层,
\(f:F\to G\) 是一个 \(S\)-态射.
则 \(f\) 是单态射, 分别地, 满态射, 分别地, 同构,
当且仅当它在 \(S\) 上局部具有相应性质.}

\medskip
\noindent\SGARefTarget{sga3:iv:remark:4:5:9}\textbf{注 4.5.9.}
\SGARefLink{sga3:iv:proposition:4:5:7}{此命题} 的证明表明,
其关于单态射 (分别地, 同构) 的断言,
即使只假设 \(F\) 是一个分离预层 (分别地, 层),
而 \(G\) 是任意预层 (分别地, 分离预层), 也仍成立.

% Translation coverage: entire assigned source component, without omissions.
% Statistics: source lines=108; target lines=121; inline math=60; display math=2; unresolved translation TODOs=0.
